\newpage
\section{Część praktyczna}
\subsection{Architektura}

\subsection{Kolejne obszary wg DevOps}
\begin{enumerate}
    \item \textbf{Plan} (opis na narzędzi do planowania zadań, przykłady rozplanowania funkcji systemu, integracja z systemem kontroli wersji)
    \begin{enumerate}
        \item Jira
    \end{enumerate}
    \item \textbf{Code} (narzędzia użyte do wytwarzania oprogramowania, alternatywy, architektura aplikacji z uwzględnieniem 12 factor app)
    \begin{enumerate}
        \item System kontroli werjsji: GitHub
        \item IDE: VS Code, IntelliJ IDEA
        \item Technologia backend: Java, Quarkus
        \item Technologia frontend: React
    \end{enumerate}
    \item \textbf{Build} (budowanie i wersjonowanie paczek)
    \begin{enumerate}
        \item GitHub Actions – proces CI/CD zapisany w kodzie
        \item Gradle – budowanie backendu
        \item Npm/Webpack – budowanie frontendu
    \end{enumerate}
    \item \textbf{Test} (etap procesu CI/CD wbudowany w GitHub Actions)
    \begin{enumerate}
        \item GitHub Actions
        \item JUnit
    \end{enumerate}
    \item \textbf{Release} (wydawanie paczek releasowych, zapisywanie kolejnych iteracji projektów)
    \begin{enumerate}
        \item Docker + DockerHub – publiczne repozytorium dla kontenerów
        \item Helm + ArtifactHub – publiczne repozytorium dla schematów Helm
    \end{enumerate}
    \item \textbf{Deploy} (infrastruktura wspierająca, klaster do uruchomienia kontenerów)
    \begin{enumerate}
        \item Terraform – tworzenie zasobów na AWS wg. koncepcji Infrastructure-as-a-Code
        \item Ansible – konfiguracja serwerów z użyciem Configuration Management
        \item Kubernetes – środowisko uruchomieniowe aplikacji
    \end{enumerate}
    \item \textbf{Operate} (środowisko chmurowe, serwery, usługi wspierające)
    \begin{enumerate}
        \item Kubernetes
            \begin{enumerate}
            \item kind
            \item minikube
        \end{enumerate}
        \item AWS Cloud Service
        \begin{enumerate}
            \item RDS – bazy danych
            \item S3 – zapis plików
            \item Opensearch – zapis dokumentów
            \item (inne usługi – rozwinąć, lista może się zmienić w trakcie implementacji)
        \end{enumerate}
    \end{enumerate}
    \item \textbf{Monitor + Feedback} (monitorowanie systemu, aplikacji, dzienniki logów, ślady aplikacji)
    \begin{enumerate}
        \item Grafana – wizualizacja metryk
        \item Prometheus – zbieranie metryk
        \item Elasticsearch/Opensearch – zbieranie logów
        \item Jaeger lub Zipkin – zbieranie trace’ów
        \item (OpenTelemetry) – standard komunikacyjny do powyższych technologii
        \item OpsGenie – alerting
    \end{enumerate}
\end{enumerate}